\documentclass[10pt]{letter}
\usepackage{url}
\usepackage[letterpaper]{geometry}

\oddsidemargin=.0in
\evensidemargin=.0in
\textwidth=5.9in
\topmargin=-1in
\textheight=10in
%\oddsidemargin=.2in
%\evensidemargin=.2in
%\textwidth=5.9in
%\topmargin=-.5in
%\textheight=9in


\name{Xiaoyu Wang\\
School of Computer Science and Technology\\
Soochow University\\
Suzhou, Jiangsu 215006, China\\
xywang21@suda.edu.cn}

\signature{Xiaoyu Wang\\
Corresponding Author}

\date{\today}

\begin{document}


% BODY OF ALL SUBLIMINAL LETTERS *******************************************

\def \body
{
We would like to submit our manuscript entitled``DARS: Task Demand-Attentive Reinforcement Scheduling for Wireless-Powered
Multi-UAV Mobile Edge Computing'' to \textit{IEEE Internet of Things Journal} for consideration of possible publication.
This paper is the result of our latest research work. It hasnot been published or accepted for publication before this submission. 
All authors have approvedthe manuscript and agree with its submission to \textit{IEEE Internet of Things Journal}. 


This manuscript investigates sustainable computation offloading for IoT
devices in a wireless-powered multi-UAV mobile edge computing system. The
problem is challenging because UAV movement, user association, task offloading,
computation and bandwidth allocation, and wireless charging are tightly
coupled, while task arrivals and spatial demand hotspots vary over time. We
therefore develop DARS, a demand-attentive reinforcement scheduling framework
that coordinates service and energy replenishment to reduce task completion
latency while maintaining long-term UAV operation.

The main contributions of this work are summarized as follows:
\begin{enumerate}
\item We introduce DARS as a cascaded scheduler that assigns long-horizon UAV
mobility and offloading control, feasibility-aware user association, and
instantaneous resource assignment to separate decision stages. This structure
keeps the learned continuous action space compact while preserving efficient
per-slot coordination.

\item We enhance policy learning with a task-demand attention encoder that
summarizes the user-demand configuration independently of user ordering. A
dual-priority replay strategy ranks samples using both value-estimation error
and reward variation and, together with multi-step returns, focuses training on
abrupt demand changes and energy-critical experiences.

\item For the per-slot computation and bandwidth subproblems, we prove strict
convexity and obtain square-root allocation rules from the KKT conditions.
These analytical solutions remove high-dimensional instantaneous resource
variables from the reinforcement-learning action space.

\item Comprehensive comparisons with representative DRL baselines, including
ESAC-PER, SAC, TD3, and MADDPG, show consistent gains in cumulative reward,
SLRR, and task completion latency across different user densities and hotspot
intensities. Ablation studies further verify the performance and stability
benefits of the attention encoder and dual-priority replay strategy.
\end{enumerate}

The topic is well aligned with the scope of the \textit{IEEE Internet of Things
Journal}, as it addresses UAV-assisted edge intelligence, computation
offloading for IoT devices, wireless energy replenishment, and adaptive
resource management in dynamic IoT environments.

This manuscript is original, has not been published previously, and is not
under consideration for publication elsewhere. All authors have approved the
manuscript and consent to its submission. The authors declare no conflict of
interest regarding this work.

Thank you for your time and consideration. Please contact me if any additional
information is required.
}

% LETTERS BEGIN HERE ******************************************************

\begin{letter}{Editorial Board\\
               IEEE Internet of Things Journal}
\opening{Dear Editor-in-Chief:} \body
\closing{Sincerely,}
\end{letter}

\end{document}
